\section{The input--output problems}\label{sec:setup}

In this section, we specify the coefficient space and local
root map, choose the output metric, and define conditioning as a derivative
norm. We then identify the matrix perturbation problem and the probability law
used to select a root.

\subsection{Coefficient space and the local root map}
Let $\mathcal H$ be an $m$-dimensional space of complex polynomials on an open
set $U\subseteq\C^d$, with a specified Hermitian inner product and orthonormal
basis $\phi_1,\ldots,\phi_m$. Write
$\phi(z)=(\phi_1(z),\ldots,\phi_m(z))^T$. 
The input is a coefficient array $c=(c_{ik})\in\C^{d\times m}$, representing
\[
 P_c=(p_1,\ldots,p_d)^T,\qquad
 p_i(z)=\sum_{k=1}^m c_{ik}\phi_k(z).
\]
We allow perturbations of every coefficient and measure perturbations using the Frobenius norm:
\[
 \|\delta c\|_F^2=\sum_{i,k}|\delta c_{ik}|^2.
\]
Thus a change of orthonormal coefficient basis changes coordinates, but not the
input norm. A nonunitary change of polynomial basis generally changes the
problem unless its Gram matrix is included in the norm.

Write $J_c(z)=(\partial p_i/\partial z_j)_{i,j}$ for the Jacobian matrix.
A root $z$ of $P_c$ is \emph{simple} if $J_c(z)$ is invertible. We will work
with the incidence set
\[
 \mathcal V=\{(c,z)\in\C^{d\times m}\times U:
                  P_c(z)=0,\ \det J_c(z)\ne0\}.
\]
At each $(c,z)\in\mathcal V$, the implicit function theorem gives a unique local
root branch 
\begin{equation}\label{eq:setup-root-branch}
   \mathcal S_{c,z}\colon c'\mapsto z(c'), \text{ with } z(c)=z.
\end{equation}   
This local map is the input-output map under study.
Its derivative at $c$ is
\begin{equation}\label{eq:setup-root-derivative}
 D\mathcal S_{c,z}[\delta c]=-J_c(z)^{-1}E_z[\delta c],
 \qquad E_z[\delta c]=\delta c\,\phi(z).
\end{equation}
Here $E_z$ evaluates the perturbed equations at the original root:
$E_z[\delta c]=\delta P(z)$. Derivatives of the root branch and other
holomorphic maps are complex-linear.

\subsection{Choosing an output metric and defining conditioning}

To quantify a change in the output, we must choose how to measure root
displacements. Fix a smooth matrix-valued function
\[
 G:U\longrightarrow\C^{d\times d},\qquad
 G(z)=G(z)^*,\qquad v^*G(z)v>0\quad(v\ne0).
\]
Thus $G$ assigns a positive definite Hermitian matrix to every possible root
location. Identifying the tangent space at
$z\in U$ with $\C^d$, define
\[
 \langle u,v\rangle_{G,z}=u^*G(z)v,\qquad
 \|v\|_{G,z}=\sqrt{v^*G(z)v}.
\]
This field of inner products thus defines a Hermitian metric on $U$. The constant
choice $G(z)=I_d$ gives ordinary Euclidean displacement. Later we will also
use a particular nonconstant choice $G=H$, determined by the coefficient
space and defined in Section~\ref{sec:setup-kernel-metric}.

\begin{definition}
For a simple system--root pair $(c,z)$, the \emph{absolute coefficient
condition number} with this output metric is defined as 
\begin{equation}\label{eq:setup-condition-definition}
 \kappa_G(c,z)
 :=\sup_{\delta c\ne0}
    \frac{\|D\mathcal S_{c,z}[\delta c]\|_{G,z}}{\|\delta c\|_F}.
\end{equation}
\end{definition}

The squared norm of scalar polynomial evaluation at $z$ is
\[
 S(z)=\|\phi(z)\|_2^2.
\]
Thus $S(z)=0$ if and only if $z$ is a common root of all the basis polynomials
$\phi_1,\ldots,\phi_m$, or equivalently of every polynomial in $\mathcal H$.
For a positive definite Hermitian matrix $G(z)$, the notation $G(z)^{1/2}$
means its unique positive definite Hermitian square root.

The next proposition specializes the standard condition-as-derivative
framework \cite[Chapter~14]{burgisser2013condition} to our coefficient space.
We give the short proof to make the evaluation factor and the choice of
output metric explicit.

\begin{proposition}\label{prop:setup-condition}
Fix an output metric $G$ as above and a simple system--root pair $(c,z)$ with
$S(z)>0$. The condition number in \eqref{eq:setup-condition-definition} is
\begin{equation}\label{eq:setup-condition}
 \kappa_G(c,z)=\sqrt{S(z)}\,
       \|G(z)^{1/2}J_c(z)^{-1}\|_2.
\end{equation}
For $G(z)=I_d$, we write $\kappa_{\mathrm{coeff}}:=\kappa_{I_d}$, so
$\kappa_{\mathrm{coeff}}(c,z)=\sqrt{S(z)}\,\|J_c(z)^{-1}\|_2$.
For a fixed holomorphic scalar output $g:U\to\C$, with its usual absolute-value
metric, the corresponding condition number is
\begin{equation}\label{eq:setup-scalar-condition}
 \kappa_{\mathrm{coeff}}^{g}(c,z)
   =\sqrt S\,\|Dg(z)J^{-1}\|_2.
\end{equation}
\end{proposition}
\begin{proof}
Suppress $c,z$ temporarily and write $J^{-*}=(J^{-1})^*$. The evaluation map
satisfies $E_zE_z^*=S(z)I_d$. Therefore
\[
 (G^{1/2}J^{-1}E_z)(G^{1/2}J^{-1}E_z)^*
   =S\,G^{1/2}J^{-1}J^{-*}G^{1/2}.
\]
Taking operator norms proves \eqref{eq:setup-condition} and the chain rule gives
\eqref{eq:setup-scalar-condition}.
\end{proof}

The assumption $S(z)>0$ is not needed for Proposition~\ref{prop:setup-condition}:
if $S(z)=0$, then $E_z=0$, so the root derivative and the displayed condition
numbers are zero, and the formulas still hold. Positivity will be important
later when defining the kernel metric through $\log S$ and normalizing the
Jacobian by $S^{-1/2}$.

The distinction between coefficient and residual perturbations is essential.
The norm $\|J^{-1}\|_2$ measures sensitivity to perturbing the equation values
$P(z)$. Perturbing the coefficient array first evaluates a polynomial at $z$,
which contributes the factor $\sqrt S$.

For nonzero input and a nonzero Euclidean output $z$, the normwise relative condition is
\begin{equation}\label{eq:setup-relative}
 \kappa_{\mathrm{coeff,rel}}(c,z)
 =\kappa_{\mathrm{coeff}}(c,z)\frac{\|c\|_F}{\|z\|_2}.
\end{equation}
For the scalar output $g(z)\ne0$, replace the denominator by $|g(z)|$ and use
$\kappa_{\mathrm{coeff}}^g$. We do not use the phrase ``relative condition''
without specifying these input and output scales.

\subsection{The kernel metric}\label{sec:setup-kernel-metric}

We now specify the second output metric used in this paper. The coefficient
space determines the reproducing kernel
\[
 K(z,w)=\sum_{k=1}^m\phi_k(z)\overline{\phi_k(w)},\qquad S(z)=K(z,z).
\]
Throughout the general local theory we assume $S>0$ and
\begin{equation}\label{eq:setup-kernel-metric}
 H_{jk}(z)=\partial_{\bar z_j}\partial_{z_k}\log S(z),
 \qquad H(z)>0\quad(z\in U).
\end{equation}
The derivatives are taken with $z$ and $\bar z$ as independent variables
(Wirtinger derivatives). The order of the indices is chosen so that a tangent
vector $v$ has squared length $v^*H(z)v$. Taking $G=H$ in the preceding
subsection defines the kernel output metric and its condition number $\kappa_H$.

There is a geometric interpretation of this formula. The notation $[\phi(z)]$
means the complex line spanned by the nonzero vector $\phi(z)$, a point of
complex projective space. The Fubini--Study metric is the standard Hermitian
metric on that space of lines. Its \emph{pullback} by $z\mapsto[\phi(z)]$
measures a displacement at $z$ by the displacement it induces in this line;
the resulting matrix is $H(z)$. Positivity means that the line changes to
first order in every nonzero tangent direction. We use the normalization
whose expression on $\mathbb P^1$ is $|\dd z|^2/(1+|z|^2)^2$.
The explicit formula \eqref{eq:setup-kernel-metric} suffices to use the metric
without this projective interpretation.

Constructing the root metric from the coefficient inner product is also
central to the sparse-system framework of Malajovich--Rojas
\cite[Sections~1--2]{malajovich2004high}; input normalizations must be
matched before comparing condition numbers.

Define the normalized Jacobian
\begin{equation}\label{eq:setup-normalized}
 A=S^{-1/2}JH^{-1/2}.
\end{equation}
Proposition~\ref{prop:setup-condition} gives the exact identity
\begin{equation}\label{eq:setup-kernel-condition}
 \kappa_H(c,z)=\|A^{-1}\|_2.
\end{equation}
This is a change of output metric, with the coefficient norm held fixed.
The Euclidean and kernel-metric condition numbers answer different questions;
Section~\ref{sec:box} relates them in the box model.

\begin{example}[Two natural metrics in one variable]\label{ex:setup-metrics}
For the monomial coefficient norm on polynomials of degree at most $n$,
$S(z)=\sum_{k=0}^n|z|^{2k}$. With $r=|z|^2$,
\[
 H(z)=h_n(r)=\left(r\frac{k_n'(r)}{k_n(r)}\right)',
 \qquad k_n(r)=\sum_{k=0}^n r^k.
\]
For the different coefficient norm making
$\sqrt{\binom nk}z^k$ orthonormal, $S(z)=(1+|z|^2)^n$ and
$H(z)=n/(1+|z|^2)^2$. The latter is the complex Kostlan model after Gaussian
coefficients are introduced. The universal local law below applies to both,
but their kernel metrics and their coefficient norms differ.
\end{example}

For comparison with the classical projective theory, consider $d$ homogeneous
equations of a common degree $n$ in $d+1$ variables with the
Bombieri--Weyl inner product. This makes the monomials
$\sqrt{n!/\alpha!}\,X^\alpha$, $|\alpha|=n$, orthonormal, where
$\alpha!=\alpha_0!\cdots\alpha_d!$.
In the affine chart $X_0=1$, their kernel satisfies
$S(z)=(1+\|z\|_2^2)^n$, so $H=nG_{\mathrm{FS}}$, with
$G_{\mathrm{FS}}$ the standard projective metric in this chart.
The classical degree-normalized projective condition number is therefore
\begin{equation}\label{eq:setup-classical-mu}
 \mu(c,[1:z])=\|c\|_F\,\kappa_H(c,z).
\end{equation}
Indeed, at $z=0$ the normalized projective Jacobian is $J/\sqrt n$;
unitary invariance gives the identity in every chart
\cite{shub1993geometric,burgisser2013condition}.
Thus $\mu$ includes both an input-norm factor and a degree normalization;
the coefficient derivative measured in $G_{\mathrm{FS}}$ has norm
$\kappa_H/\sqrt n$. Neither should be identified with the affine relative
condition in \eqref{eq:setup-relative}.

\subsection{The matrix eigenvalue problem and its perturbations}

For an eigenvalue $\lambda$ of algebraic multiplicity one of
$M\in\C^{N\times N}$, choose nonzero
vectors $r,w$ such that $Mr=\lambda r$ and $w^*M=\lambda w^*$.
The local eigenvalue map $\mathcal E_{M,\lambda}$ sends a nearby matrix to
its unique eigenvalue near $\lambda$. Its derivative at $M$ is
\[
 D\mathcal E_{M,\lambda}[\delta M]=\frac{w^*\delta M r}{w^*r}.
\]
With arbitrary matrix perturbations measured in spectral norm and scalar
output measured in absolute value, its condition is
\begin{equation}\label{eq:setup-overlap}
 \kappa_{\mathrm{eig}}(M,\lambda)
 =\Lambda(M,\lambda):=\frac{\|w\|_2\|r\|_2}{|w^*r|}\ge1.
\end{equation}
The derivative norm is at most the displayed overlap, and a rank-one
perturbation attains it; $\Lambda\ge1$ follows from Cauchy--Schwarz.
The same absolute condition is obtained using Frobenius norm,
but its relative version would then use $\|M\|_F$. Here we consistently use
\begin{equation}\label{eq:setup-eigen-relative}
 \kappa_{\mathrm{eig,rel}}(M,\lambda)
 =\Lambda(M,\lambda)\frac{\|M\|_2}{|\lambda|}
 \quad(\lambda\ne0).
\end{equation}

Suppose a differentiable algebraic construction $\mathcal D:c\mapsto M(c)$
represents $g(z)$ as an eigenvalue. On compatible local branches,
\begin{equation}\label{eq:setup-composition}
 \mathcal E\circ\mathcal D=g\circ\mathcal S,
 \qquad D\mathcal E\,D\mathcal D=Dg\,D\mathcal S.
\end{equation}
Consequently coefficient-induced matrix perturbations give exactly the
original derivative. The comparison in this paper instead allows arbitrary
perturbations in the intermediate matrix space. It describes the sensitivity
of that larger problem. Errors in constructing $M$, implementing an eigensolver,
and recovering other coordinates require additional analysis.

One important construction uses the quotient algebra of the roots. Suppose
the ideal $I=\langle p_1,\ldots,p_d\rangle$ has exactly $N$ distinct complex
zeros $z_1,\ldots,z_N$ and consists of all polynomials vanishing at these
points (equivalently, $I$ is radical). In
$\mathcal A=\C[x_1,\ldots,x_d]/I$, two polynomials represent the same element
precisely when their values at all roots agree. Evaluation identifies
$\mathcal A$ with the $N$-dimensional space of functions on this root set.
Choose a basis $\B=(b_1,\ldots,b_N)$ of polynomial representatives.
For a polynomial $g$, multiplication by its class is a linear map on
$\mathcal A$; denote its matrix, acting on coefficient columns, by $M_g$.
The evaluation matrix $V=(b_k(z_j))_{k,j}$ is invertible, and evaluation gives
\[
 M_g=V^{-T}\diag(g(z_1),\ldots,g(z_N))V^T.
\]
If $g$ \emph{separates} the roots, meaning that $g(z_1),\ldots,g(z_N)$ are
pairwise distinct, its overlap at $z_j$ is therefore
\begin{equation}\label{eq:setup-interpolation}
 \Lambda_j=\|V^{-T}e_j\|_2\|Ve_j\|_2,
 \qquad
 \Lambda_j^2=
 \frac{\|Ve_j\|_2^2}
 {\dist(Ve_j,\spanop\{Ve_k:k\ne j\})^2}.
\end{equation}
The last identity follows by orthogonally projecting one evaluation column
away from the others. Thus matrix conditioning becomes a question about
interpolation in a specified basis norm. This deterministic interpretation
also underlies the normal-form analysis in \cite{graf2026numerical}.

\subsection{Gaussian input and a uniformly selected root}\label{sec:setup-probability}

We now make the input random: all $dm$ coefficient coordinates are independent
$\CN(0,1)$ variables, with density $\pi^{-1}e^{-|a|^2}$ on $\C$.
Equivalently, the real and imaginary parts have variance $1/2$.
The components of $P$ are independent and share the kernel $K$.

Assume that almost surely $P$ has exactly $N\ge1$ zeros in $U$, all simple,
where $N$ is deterministic. First sample $P$ with this Gaussian law, then select
one of its roots uniformly. For a nonnegative measurable function $F(P,z)$,
called a \emph{mark}, define
\begin{equation}\label{eq:setup-selected-root}
 \Estar F(P,Z)=\frac1N\E\sum_{P(z)=0,\ z\in U}F(P,z).
\end{equation}
This determines a probability law $\Pstar$ on system--root pairs. In particular,
the marginal law of $P$ is its original Gaussian law. This last fact uses the
deterministic root count. No independence between the roots of one system is
assumed.

Sampling a root this way is different from conditioning the Gaussian input on
$P(z)=0$ at a prescribed point. The former also weights the conditional law
by $|\det J(z)|^2$, as Section~\ref{sec:local-law} proves. We use conditional
expectations given $Z=z$ in the regular conditional sense supplied by that
formula, rather than treating a zero-probability event as an ordinary event.

For example, for an overlap threshold $t$,
\[
 \Pstar(\Lambda>t)=\frac1N\E\#\{j:\Lambda_j>t\}.
\]
A tail bound therefore controls the expected fraction of problematic roots.
It does not give the same bound for the worst root of a system. Likewise
$\Estar\log X$ is a logarithmic average, not a median or a high-probability
bound for $X$.
